\documentclass[../../main.tex]{subfiles}% 注意这里的文件路径不能用 ./main.tex ,否则用latexmk编译子文件会报错
\graphicspath{{\subfix{./image/}}{\subfix{../../image/}}} % 指定图片引用路径,后续可以直接使用图片文件名
% 如果图片存放在上级目录,则使用 ../../image/ 作为备用路径

% 例如:
% \begin{figure}[H]
% \centering
% \includegraphics[width=0.3\textwidth]{图.png}
% \caption{}
% \label{figure:图}
% \end{figure}
% 注意:上述\label{}一定要放在\caption{}之后,否则引用图片序号会只会显示??.

\begin{document}

\section{混合问题}

\subsection{分离变量法}

\begin{definition}
称如下问题为\textbf{一维波动方程的混合问题}：
\begin{align}
\begin{cases}
u_{tt}-a^2u_{xx}=f, & (x,t)\in(0,l)\times(0,+\infty), \\
u(x,0)=\varphi(x), & x\in[0,l], \\
u_t(x,0)=\psi(x), & x\in[0,l], \\
u(0,t)=g_1(t),\quad u(l,t)=g_2(t), & t\in[0,+\infty).
\end{cases} \label{eq:wave-mixed}
\end{align}
其中$f,\varphi ,\psi,g_1,g_2$都是已知函数.
\end{definition}

\begin{theorem}[一维波动方程的混合问题齐次方程、齐次边值情形的解]\label{theorem:一维波动方程的混合问题齐次方程、齐次边值情形的解}
设 $\varphi(x)\in C^3[0,l]$，$\psi(x)\in C^2[0,l]$ 满足相容性条件
\begin{align}
\varphi(0)=\varphi(l)=\varphi''(0)=\varphi''(l)=\psi(0)=\psi(l)=0. \label{eq:comp-mixed}
\end{align}
则混合问题
\begin{align}
\begin{cases}
u_{tt}-a^2u_{xx}=0, & (x,t)\in(0,l)\times(0,+\infty), \\
u(x,0)=\varphi(x), & x\in[0,l], \\
u_t(x,0)=\psi(x), & x\in[0,l], \\
u(0,t)=u(l,t)=0, & t\in[0,+\infty)
\end{cases} \label{eq:wave-mixed-hom}
\end{align}
的解$u\in C^2(\overline{Q})$,其中 $Q=(0,l)\times(0,+\infty)$.且$u$可表示成
\begin{align}
u(x,t)=\sum_{n=1}^{\infty}\left(A_n\cos\frac{na\pi}{l}t+B_n\sin\frac{na\pi}{l}t\right)\sin\frac{n\pi}{l}x, \label{eq:wave-sep-sol}
\end{align}
其中
\begin{gather}
A_n=\frac{2}{l}\int_0^l\varphi(x)\sin\frac{n\pi}{l}x\,\mathrm{d}x,\quad n=1,2,\cdots, \label{eq:wave-A} \\
B_n=\frac{2}{na\pi}\int_0^l\psi(x)\sin\frac{n\pi}{l}x\,\mathrm{d}x,\quad n=1,2,\cdots. \label{eq:wave-B}
\end{gather}
\end{theorem}
\begin{remark}
相容性条件 \eqref{eq:comp-mixed} 的导出与半无界问题的情形类似，它保证了解在角点 $(0,0)$ 和 $(l,0)$ 处的光滑性.
\end{remark}
\begin{proof}
采用分离变量法。设非零解具有形式
\begin{align*}
u(x,t)=X(x)T(t),
\end{align*}
其中 $X\not\equiv 0$，$T\not\equiv 0$. 代入方程 $u_{tt}-a^2u_{xx}=0$，得
\begin{align*}
T''(t)X(x)-a^2X''(x)T(t)=0,\qquad (x,t)\in(0,l)\times(0,+\infty),
\end{align*}
即
\begin{align*}
\frac{T''(t)}{a^2T(t)}=\frac{X''(x)}{X(x)}.
\end{align*}
上式左端仅与 $t$ 有关，右端仅与 $x$ 有关，故必为常数，记为 $-\lambda$，于是
\begin{align*}
T''(t)+a^2\lambda T(t)=0,\qquad X''(x)+\lambda X(x)=0.
\end{align*}
将 $u=X(x)T(t)$ 代入边值条件 $u(0,t)=u(l,t)=0$，得
\begin{align*}
X(0)T(t)=X(l)T(t)=0,\qquad t>0.
\end{align*}
因 $T(t)\not\equiv 0$，故 $X(0)=X(l)=0$. 从而得到特征值问题
\begin{align*}
\begin{cases}
X''(x)+\lambda X(x)=0,\quad 0<x<l,\\
X(0)=X(l)=0.
\end{cases}
\end{align*}
求解该问题。若 $\lambda\leqslant 0$，则结合边值条件只能得到零解，与 $X\not\equiv 0$ 矛盾，故 $\lambda>0$. 令 $\mu=\sqrt{\lambda}$，则通解为
\begin{align*}
X(x)=C_1\cos \mu x+C_2\sin \mu x.
\end{align*}
由 $X(0)=0$ 得 $C_1=0$；由 $X(l)=0$ 得 $C_2\sin \mu l=0$. 因 $C_2\neq 0$，所以 $\sin \mu l=0$，即 $\mu l=n\pi$，$n=1,2,\cdots$. 故特征值和对应的特征函数为
\begin{align*}
\lambda_n=\left(\frac{n\pi}{l}\right)^2,\qquad X_n(x)=\sin\frac{n\pi}{l}x,\quad n=1,2,\cdots.
\end{align*}
对于每个 $n$，方程 $T''+a^2\lambda_n T=0$ 的通解为
\begin{align*}
T_n(t)=A_n\cos\frac{na\pi}{l}t+B_n\sin\frac{na\pi}{l}t,
\end{align*}
其中 $A_n,B_n$ 为任意常数。于是得到满足方程和边值条件的一系列特解
\begin{align*}
u_n(x,t)=\left(A_n\cos\frac{na\pi}{l}t+B_n\sin\frac{na\pi}{l}t\right)\sin\frac{n\pi}{l}x,\quad n=1,2,\cdots.
\end{align*}
由叠加原理，形式上构造
\begin{align*}
u(x,t)=\sum_{n=1}^{\infty}u_n(x,t)
=\sum_{n=1}^{\infty}\left(A_n\cos\frac{na\pi}{l}t+B_n\sin\frac{na\pi}{l}t\right)\sin\frac{n\pi}{l}x.
\end{align*}
为使 $u$ 满足初始条件，令
\begin{align}\label{eq::jioj33jr2333333}
u(x,0)=\varphi(x)=\sum_{n=1}^{\infty}A_n\sin\frac{n\pi}{l}x,\qquad
u_t(x,0)=\psi(x)=\sum_{n=1}^{\infty}\frac{na\pi}{l}B_n\sin\frac{n\pi}{l}x.
\end{align}
根据特征函数系 $\left\{ \sin \left( \frac{n\pi x}{l} \right) \right\} $ 在 $[0,l]$ 上的完备性，$\varphi$ 和 $\psi$ 可展开为 Fourier 正弦级数
\begin{align*}
\varphi(x)=\sum_{n=1}^{\infty}\varphi_n\sin\frac{n\pi}{l}x,\qquad
\psi(x)=\sum_{n=1}^{\infty}\psi_n\sin\frac{n\pi}{l}x,
\end{align*}
上式两边同乘$\sin\frac{n\pi}{l}x(n\geqslant 1)$再对$x$积分得
\begin{align*}
\varphi_n=\frac{2}{l}\int_0^l\varphi(x)\sin\frac{n\pi}{l}x\,\mathrm{d}x,\qquad
\psi_n=\frac{2}{l}\int_0^l\psi(x)\sin\frac{n\pi}{l}x\,\mathrm{d}x.
\end{align*}
代入\eqref{eq::jioj33jr2333333}式得
\begin{align*}
A_n=\varphi_n=\frac{2}{l}\int_0^l\varphi(x)\sin\frac{n\pi}{l}x\,\mathrm{d}x,\qquad
B_n=\frac{l}{na\pi}\psi_n=\frac{2}{na\pi}\int_0^l\psi(x)\sin\frac{n\pi}{l}x\,\mathrm{d}x.
\end{align*}
至此得到形式解。

下面证明该级数确实属于 $C^2(\overline{Q})$ 且满足方程和初边值条件。由相容性条件 $\varphi(0)=\varphi(l)=\varphi''(0)=\varphi''(l)=0$，对 $A_n$ 连续三次分部积分，得
\begin{align*}
A_n=-\frac{l^3}{\pi^3 n^3}a_n,\qquad a_n=\frac{2}{l}\int_0^l\varphi'''(x)\sin\frac{n\pi}{l}x\,\mathrm{d}x.
\end{align*}
同理，利用 $\psi(0)=\psi(l)=0$，两次分部积分得
\begin{align*}
B_n=-\frac{l^3}{a\pi^3 n^3}b_n,\qquad b_n=\frac{2}{l}\int_0^l\psi''(x)\sin\frac{n\pi}{l}x\,\mathrm{d}x.
\end{align*}
记 $u_n(x,t)$ 为级数的通项，则存在常数 $C_1,C_2,C_3$，使得在 $\overline{Q}=[0,l]\times[0,+\infty)$ 上
\begin{align*}
|u_n(x,t)|\leqslant \frac{C_1}{n^3},\qquad |Du_n(x,t)|\leqslant \frac{C_2}{n^2},\qquad |D^2u_n(x,t)|\leqslant a_n^2+b_n^2+\frac{C_3}{n^2}.
\end{align*}
由 Bessel 不等式，$\sum a_n^2$ 与 $\sum b_n^2$ 收敛，故级数 $\sum u_n$，$\sum Du_n$，$\sum D^2u_n$ 均在 $\overline{Q}$ 上一致收敛。因此 $u\in C^2(\overline{Q})$，且可逐项求导两次。直接代入波动方程，并利用每一项 $u_n$ 均满足方程，可得 $u$ 满足方程；代入边值条件，因每一项在边界为零，故 $u$ 满足齐次边值；代入初始条件，由 Fourier 系数的确定，$u(x,0)=\varphi(x)$，$u_t(x,0)=\psi(x)$. 综上，该级数即为混合问题的解。

由相容性条件 $\varphi(0)=\varphi(l)=\varphi''(0)=\varphi''(l)=0$，利用分部积分得
\begin{align*}
A_n=-\frac{l^3}{\pi^3n^3}a_n,\qquad a_n=\frac{2}{l}\int_0^l\varphi'''(x)\sin\frac{n\pi}{l}x\,\mathrm{d}x.
\end{align*}
同理，
\begin{align*}
B_n=-\frac{l^3}{a\pi^3n^3}b_n,\qquad b_n=\frac{2}{l}\int_0^l\psi''(x)\sin\frac{n\pi}{l}x\,\mathrm{d}x.
\end{align*}
于是存在常数 $C_1,C_2,C_3$，使得在 $\overline{Q}=[0,l]\times[0,+\infty)$ 上
\begin{align*}
|u_n(x,t)|\leqslant \frac{C_1}{n^3},\qquad |Du_n(x,t)|\leqslant \frac{C_2}{n^2},\qquad |D^2u_n(x,t)|\leqslant a_n^2+b_n^2+\frac{C_3}{n^2}.
\end{align*}
由 Bessel 不等式，$\sum a_n^2$，$\sum b_n^2$ 收敛，故级数 $\sum u_n$，$\sum Du_n$，$\sum D^2u_n$ 在 $\overline{Q}$ 上一致收敛. 因此 $u\in C^2(\overline{Q})$，且可逐项求导两次，直接代入方程和初边值条件即验证其为解.
\end{proof}

\begin{theorem}[一维波动方程的混合问题非齐次方程、齐次边值情形的解]\label{theorem:一维波动方程的混合问题非齐次方程、齐次边值情形的解}
设 $\varphi\in C^3[0,l]$，$\psi\in C^2[0,l]$，$f\in C^2([0,l]\times[0,+\infty))$ 满足相容性条件
\begin{gather}
\varphi(0)=\varphi(l)=\psi(0)=\psi(l)=0, \,\,
a^2\varphi''(0)+f(0,0)=0,\,\,a^2\varphi''(l)+f(l,0)=0. \label{eq:comp-mixed-nonhom}
\end{gather}
则混合问题
\begin{align}
\begin{cases}
u_{tt}-a^2u_{xx}=f, & (x,t)\in(0,l)\times(0,+\infty), \\
u(x,0)=\varphi(x), & x\in[0,l], \\
u_t(x,0)=\psi(x), & x\in[0,l], \\
u(0,t)=u(l,t)=0, & t\in[0,+\infty)
\end{cases} \label{eq:wave-mixed-nonhom}
\end{align}
的解$u\in C^2(\overline{Q})$,其中 $Q=(0,l)\times(0,+\infty)$.且可表示成
\begin{align}
u(x,t)=\sum_{n=1}^{\infty}T_n(t)\sin\frac{n\pi}{l}x, \label{eq:wave-u-series}
\end{align}
其中
\begin{align}
T_n(t)=\varphi_n\cos\frac{n\pi a}{l}t+\frac{l}{n\pi a}\psi_n\sin\frac{n\pi a}{l}t 
+\frac{l}{n\pi a}\int_0^t f_n(\tau)\sin\frac{n\pi a}{l}(t-\tau)\,\mathrm{d}\tau, \label{eq:wave-Tn}
\end{align}
而
\begin{align*}
\varphi_n=\frac{2}{l}\int_0^l\varphi(x)\sin\frac{n\pi}{l}x\,\mathrm{d}x,\quad
\psi_n=\frac{2}{l}\int_0^l\psi(x)\sin\frac{n\pi}{l}x\,\mathrm{d}x, \quad
f_n(t)=\frac{2}{l}\int_0^l f(x,t)\sin\frac{n\pi}{l}x\,\mathrm{d}x. 
\end{align*}
\end{theorem}
\begin{remark}
由 \eqref{eq:wave-Tn} 可知，非齐次项 $f$ 的解可看作每个模态 $n$ 的受迫振动，其中 $f_n(t)$ 是外力在对应模态上的投影，$\frac{l}{n\pi a}\int_0^t f_n(\tau)\sin\frac{n\pi a}{l}(t-\tau)\,\mathrm{d}\tau$ 是 Duhamel 积分.
\end{remark}
\begin{proof}
将 $u$，$f$，$\varphi$，$\psi$ 按特征函数系 $\{\sin(n\pi x/l)\}$ 展开，代入方程得
\begin{align*}
T_n''(t)+\left(\frac{n\pi a}{l}\right)^2T_n(t)=f_n(t),\qquad T_n(0)=\varphi_n,\quad T_n'(0)=\psi_n.
\end{align*}
解此常微分方程初值问题即得 \eqref{eq:wave-Tn}. 由相容性条件 \eqref{eq:comp-mixed-nonhom} 及分部积分可保证级数 \eqref{eq:wave-u-series} 及其一、二阶导数在 $\overline{Q}$ 上一致收敛，从而 $u\in C^2(\overline{Q})$ 且满足方程和初边值条件.
\end{proof}

\begin{theorem}[一维波动方程的混合问题非齐次方程、非齐次边值情形的解]
\label{theorem:mixed-nonhom-bc}
设 $\varphi\in C^3[0,l]$，$\psi\in C^2[0,l]$，$f\in C^2([0,l]\times[0,+\infty))$，$g_1,g_2\in C^4[0,+\infty)$ 满足相容性条件
\begin{gather}
\varphi(0)=g_1(0),\qquad \varphi(l)=g_2(0), \label{eq:comp-bc1} \\
\psi(0)=g_1'(0),\qquad \psi(l)=g_2'(0), \label{eq:comp-bc2} \\
g_1''(0)-a^2\varphi''(0)=f(0,0),\qquad g_2''(0)-a^2\varphi''(l)=f(l,0). \label{eq:comp-bc3}
\end{gather}
则混合问题
\begin{align}
\begin{cases}
u_{tt}-a^2u_{xx}=f, & (x,t)\in(0,l)\times(0,+\infty), \\
u(x,0)=\varphi(x), & x\in[0,l], \\
u_t(x,0)=\psi(x), & x\in[0,l], \\
u(0,t)=g_1(t),\quad u(l,t)=g_2(t), & t\in[0,+\infty)
\end{cases} \label{eq:mixed-nonhom-bc}
\end{align}
的解$u\in C^2(\overline{Q})$,其中$Q=(0,l)\times(0,+\infty)$,且$u$可表示成
\begin{align}
u(x,t)=h(x,t)+v(x,t), \label{eq:mixed-sol-bc}
\end{align}
其中
\begin{align}
h(x,t)=\frac{l-x}{l}g_1(t)+\frac{x}{l}g_2(t), \label{eq:h-def}
\end{align}
而 $v$ 是下列齐次边值混合问题的解：
\begin{align}
\begin{cases}
v_{tt}-a^2v_{xx}=F(x,t), & (x,t)\in Q, \\
v(x,0)=\varphi_v(x), & x\in[0,l], \\
v_t(x,0)=\psi_v(x), & x\in[0,l], \\
v(0,t)=v(l,t)=0, & t\in[0,+\infty),
\end{cases} \label{eq:v-problem}
\end{align}
其中
\begin{gather*}
\varphi_v(x)=\varphi(x)-h(x,0)=\varphi(x)-\frac{l-x}{l}g_1(0)-\frac{x}{l}g_2(0),  \\
\psi_v(x)=\psi(x)-h_t(x,0)=\psi(x)-\frac{l-x}{l}g_1'(0)-\frac{x}{l}g_2'(0),  \\
F(x,t)=f(x,t)-h_{tt}(x,t)=f(x,t)-\frac{l-x}{l}g_1''(t)-\frac{x}{l}g_2''(t).
\end{gather*}
具体地，$v$ 的级数表达式为
\begin{align}
v(x,t)=\sum_{n=1}^{\infty}T_n(t)\sin\frac{n\pi}{l}x, \label{eq:v-series}
\end{align}
其中
\begin{align}
T_n(t)=\varphi_{v,n}\cos\frac{na\pi}{l}t+\frac{l}{na\pi}\psi_{v,n}\sin\frac{na\pi}{l}t +\frac{l}{na\pi}\int_0^t F_n(\tau)\sin\frac{na\pi}{l}(t-\tau)\,\mathrm{d}\tau, \label{eq:Tn-bc}
\end{align}
而
\begin{align}
\varphi_{v,n}=\frac{2}{l}\int_0^l\varphi_v(x)\sin\frac{n\pi}{l}x\,\mathrm{d}x,\quad
\psi_{v,n}=\frac{2}{l}\int_0^l\psi_v(x)\sin\frac{n\pi}{l}x\,\mathrm{d}x,\quad
F_n(t)=\frac{2}{l}\int_0^l F(x,t)\sin\frac{n\pi}{l}x\,\mathrm{d}x. \label{eq:coeff-v}
\end{align}
\end{theorem}
\begin{remark}
对于齐次方程、齐次边值情形（即 $f\equiv 0$，$g_1\equiv g_2\equiv 0$），也可采用对称开拓法求解：将初值 $\varphi,\psi$ 奇延拓并周期延拓为 $2l$，然后利用 D'Alembert 公式得到全直线上的解，再限制到 $[0,l]$ 上，此解与级数解 \eqref{eq:wave-sep-sol} 相同. 对称开拓法对初值的光滑性要求可降低为 $\varphi\in C^2[0,l]$，$\psi\in C^1[0,l]$.
\end{remark}
\begin{remark}
当 $g_1,g_2$ 非齐次时，上述变换将非齐次边值问题化为齐次边值问题，但代价是非齐次项和初值发生了变化. 这正是处理非齐次边值问题的标准方法.
\end{remark}
\begin{proof}
作函数变换 $u=h+v$，其中 $h$ 由 \eqref{eq:h-def} 定义. 直接计算得 $h_{xx}=0$，$h_{tt}=\frac{l-x}{l}g_1''(t)+\frac{x}{l}g_2''(t)$. 于是
\begin{align*}
u_{tt}-a^2u_{xx}=(v_{tt}+h_{tt})-a^2v_{xx}=f
\Longrightarrow v_{tt}-a^2v_{xx}=f-h_{tt}=F.
\end{align*}
初值条件：
\begin{align*}
v(x,0)=u(x,0)-h(x,0)=\varphi(x)-h(x,0)=\varphi_v(x),\\
v_t(x,0)=u_t(x,0)-h_t(x,0)=\psi(x)-h_t(x,0)=\psi_v(x).
\end{align*}
边值条件：
\begin{align*}
v(0,t)=u(0,t)-h(0,t)=g_1(t)-g_1(t)=0,\\
v(l,t)=u(l,t)-h(l,t)=g_2(t)-g_2(t)=0.
\end{align*}
因此 $v$ 满足非齐次方程、齐次边值的混合问题 \eqref{eq:v-problem}. 由相容性条件 \eqref{eq:comp-bc1}\eqref{eq:comp-bc2}\eqref{eq:comp-bc3} 可推出 $\varphi_v(0)=\varphi_v(l)=\psi_v(0)=\psi_v(l)=0$ 以及
\begin{align*}
a^2\varphi_v''(0)+F(0,0)=a^2\varphi''(0)+f(0,0)-g_1''(0)=0,\\
a^2\varphi_v''(l)+F(l,0)=a^2\varphi''(l)+f(l,0)-g_2''(0)=0,
\end{align*}
即 $\varphi_v,\psi_v,F$ 满足\hyperref[theorem:一维波动方程的混合问题非齐次方程、齐次边值情形的解]{非齐次方程齐次边值定理}的相容性条件. 此外，由 $\varphi\in C^3$，$\psi\in C^2$，$f\in C^2$，$g_1,g_2\in C^4$ 可知 $\varphi_v\in C^3$，$\psi_v\in C^2$，$F\in C^2$. 故应用该定理，$v$ 可表示为级数 \eqref{eq:v-series}\eqref{eq:Tn-bc}，且 $v\in C^2(\overline{Q})$. 从而 $u=h+v\in C^2(\overline{Q})$，直接代入原问题即验证其为解.
\end{proof}



\subsection{驻波法与共振}

\begin{definition}
设混合问题 \eqref{eq:wave-mixed-hom} 的解由级数 \eqref{eq:wave-sep-sol} 给出。称其第 $n$ 项
\begin{align}
u_n(x,t)=N_n\sin\frac{n\pi}{l}x\sin\left(\frac{na\pi}{l}t+\alpha_n\right), \label{eq:standing-wave}
\end{align}
其中
\begin{align*}
N_n=\sqrt{A_n^2+B_n^2},\qquad \alpha_n=\arctan\frac{A_n}{B_n},
\end{align*}
为弦的第 $n$ 个\textbf{振动元素}（或\textbf{驻波}）。
\end{definition}

\begin{proposition}
驻波 $u_n(x,t)$ 具有以下性质：
\begin{enumerate}[(1)]
\item 频率为 $\omega_n=\dfrac{na\pi}{l}$，初相位为 $\alpha_n$；
\item 振幅为 $a_n(x)=N_n\sin\dfrac{n\pi}{l}x$，在点 $x=\dfrac{kl}{n}$（$k=0,1,\cdots,n$）处为零（节点），在点 $x=\dfrac{(2k+1)l}{2n}$（$k=0,1,\cdots,n-1$）处取极值（波腹）；
\item 每个驻波满足波动方程和齐次边值条件 $u(0,t)=u(l,t)=0$.
\end{enumerate}
\end{proposition}
\begin{remark}
分离变量法在物理学中称为\textbf{驻波法}，其物理意义是：无外力作用时，弦的振动可分解为一系列具有特定频率的驻波的叠加.
\end{remark}
\begin{remark}
在弦乐器演奏中，最低频率 $\omega_1=\dfrac{\pi}{l}\sqrt{\dfrac{T_0}{\rho_0}}$（$T_0$ 为张力，$\rho_0$ 为密度）对应的振动元素 $u_1$ 称为\textbf{基音}，其振幅 $N_1$ 通常远大于其他振动元素的振幅，决定音乐的基调；其余 $u_n$（$n\geqslant 2$）称为\textbf{泛音}，决定音色.
\end{remark}
\begin{remark}
由 $\omega_1=\dfrac{\pi}{l}\sqrt{\dfrac{T_0}{\rho_0}}$ 可知：
\begin{enumerate}[(i)]
\item 按压弦线改变有效长度 $l$ 可改变音调（长度减半则频率加倍，提高八度）；
\item 调节张力 $T_0$ 可改变音调（张力增大则音调升高）；
\item 弦的粗细（密度 $\rho_0$）影响音调（细弦音调高，粗弦音调低）.
\end{enumerate}
\end{remark}

\begin{theorem}[共振]
考虑非齐次混合问题
\begin{align}
\begin{cases}
u_{tt}-a^2u_{xx}=A(x)\sin\omega t, & (x,t)\in(0,l)\times(0,+\infty), \\
u(x,0)=0,\quad u_t(x,0)=0, & x\in[0,l], \\
u(0,t)=u(l,t)=0, & t\in[0,+\infty),
\end{cases} \label{eq:resonance}
\end{align}
其中 $\omega>0$，$A(x)\in C[0,l]$，$A(0)=A(l)=0$. 则：
\begin{enumerate}[(1)]
\item 若 $\omega\neq \omega_n=\dfrac{n\pi a}{l}$（$n=1,2,\cdots$），则解一致有界；
\item 若存在 $k\in\mathbb{N}$ 使得 $\omega=\omega_k$ 且 $a_k\neq 0$，其中
\begin{align}
a_k=\frac{2}{l}\int_0^l A(x)\sin\frac{k\pi}{l}x\,\mathrm{d}x, \label{eq:forcing-coef}
\end{align}
则解中含有随 $t$ 线性增长的项，振幅趋于无穷.
\end{enumerate}
\end{theorem}
\begin{remark}
上述定理中的现象称为\textbf{共振}：当周期外力的频率 $\omega$ 等于系统的某个固有频率 $\omega_k$ 时，若外力在该模态上的投影 $a_k\neq 0$，则对应驻波的振幅将随时间无限增大，最终导致弦的断裂.
\end{remark}
\begin{remark}
共振现象表明波动方程解不具有极大模估计（即解的最大值不能仅由边界和初值控制），这与位势方程和热方程的性质有本质区别.
\end{remark}
\begin{proof}
由 \eqref{eq:wave-u-series} 和 \eqref{eq:wave-Tn}，注意到 $\varphi=\psi=0$，得
\begin{align*}
u(x,t)=\sum_{n=1}^{\infty}\frac{a_n}{\omega_n}\sin\frac{n\pi}{l}x\int_0^t\sin\omega\tau\sin\omega_n(t-\tau)\,\mathrm{d}\tau,
\end{align*}
其中 $a_n$ 如 \eqref{eq:forcing-coef} 定义，$\omega_n=\dfrac{n\pi a}{l}$.

当 $\omega\neq \omega_n$ 时，
\begin{align*}
\int_0^t\sin\omega\tau\sin\omega_n(t-\tau)\,\mathrm{d}\tau
=\frac{\omega\sin\omega_nt-\omega_n\sin\omega t}{\omega_n(\omega^2-\omega_n^2)},
\end{align*}
故 $|u_n|$ 一致有界.

当 $\omega=\omega_k$ 时，
\begin{align*}
\int_0^t\sin\omega_k\tau\sin\omega_k(t-\tau)\,\mathrm{d}\tau
=\frac{1}{2\omega_k^2}(\sin\omega_kt-t\omega_k\cos\omega_kt).
\end{align*}
若 $a_k\neq 0$，则第 $k$ 项含有因子 $t\cos\omega_kt$，振幅随 $t\to\infty$ 线性增长.
\end{proof}



\subsection{能量不等式}

\begin{theorem}
设 $\Omega$ 是 $\mathbb{R}^n$ 上的一个有界开集，并令 $\Omega_T=\Omega\times(0,T)$。如果 $u\in C^1(\overline{\Omega}_T)\cap C^2(\Omega_T)$ 是混合问题
\begin{align}
\begin{cases}
u_{tt}-a^2\Delta u=f, & (x,t)\in\Omega\times(0,T), \\
u(x,0)=\varphi(x), & x\in\overline{\Omega}, \\
u_t(x,0)=\psi(x), & x\in\overline{\Omega}, \\
u(x,t)=0, & t\in[0,T],\ x\in\partial\Omega,
\end{cases} \label{eq:wave-mixed-energy-gen}
\end{align}
的解，则存在只依赖于 $T$ 的常数 $M$，使得对于任意 $\tau\in[0,T]$，有
\begin{align}
\int_\Omega [u_t^2(x,\tau)+a^2|Du(x,\tau)|^2]\,\mathrm{d}x
\leqslant M\left\{\int_\Omega[\psi^2(x)+a^2|D\varphi(x)|^2]\,\mathrm{d}x+\iint_{\Omega_\tau}f^2(x,t)\,\mathrm{d}x\,\mathrm{d}t\right\}, \label{eq:energy-est-wave}
\end{align}
其中 $\Omega_\tau=\Omega\times(0,\tau)$。当 $f\equiv 0$ 时，$M=1$ 且上式中的不等号由等号代替。
\end{theorem}

\begin{theorem}
设 $\Omega$ 是 $\mathbb{R}^2$ 上的一个有界开集，$\Omega_T=\Omega\times(0,T)$。如果 $u\in C^1(\overline{\Omega}_T)\cap C^2(\Omega_T)$ 是混合问题 \eqref{eq:wave-mixed-energy-gen} 的解，则存在只依赖于 $T$ 的常数 $M$，使得对于 $\tau\in[0,T]$，有
\begin{align}
&\iint_\Omega [u_t^2(x,y,\tau)+a^2(u_x^2(x,y,\tau)+u_y^2(x,y,\tau))]\,\mathrm{d}x\mathrm{d}y
\nonumber \\
&\leqslant M\Bigg\{\iint_\Omega[\psi^2(x,y)+a^2(\varphi_x^2(x,y)+\varphi_y^2(x,y))]\,\mathrm{d}x\mathrm{d}y 
+\iiint_{\Omega_\tau}f^2(x,y,t)\,\mathrm{d}x\mathrm{d}y\mathrm{d}t\Bigg\}, \label{eq:energy-est-2d}
\end{align}
其中 $\Omega_\tau=\Omega\times(0,\tau)$（见\reffig{figure:omega_tau}）。当 $f\equiv 0$ 时，$M=1$ 且上式中的不等号由等号代替。
\end{theorem}
\begin{remark}
若不记常数因子，表达式
\begin{align*}
\iint_\Omega [u_t^2(x,y,\tau)+a^2(u_x^2(x,y,\tau)+u_y^2(x,y,\tau))]\,\mathrm{d}x\mathrm{d}y
\end{align*}
表示物体 $\Omega$ 在 $\tau$ 时刻的总能量，其中第一项为动能，第二项为弹性势能。当 $f\equiv 0$ 时，总能量守恒。
\end{remark}
\begin{remark}
对不等式 \eqref{eq:energy-est-2d} 关于 $\tau$ 从 $0$ 到 $t$ 积分，可得到另一种形式的能量模估计
\begin{align}
\iiint_{\Omega_t}[u_t^2+a^2(u_x^2+u_y^2)]\,\mathrm{d}x\mathrm{d}y\mathrm{d}t
\leqslant M\left\{\iint_\Omega[\psi^2+a^2(\varphi_x^2+\varphi_y^2)]\,\mathrm{d}x\mathrm{d}y+\iiint_{\Omega_t}f^2\,\mathrm{d}x\mathrm{d}y\mathrm{d}t\right\}, \label{eq:energy-est-int}
\end{align}
其中 $M$ 只依赖于 $T$，$\Omega_t=\Omega\times(0,t)$。同样可由此得到解 $u$ 的 $L^2$ 估计。
\end{remark}
\begin{remark}
从能量模估计 \eqref{eq:energy-est-wave} 不难导出混合问题解在函数类 $C^1(\overline{\Omega}_T)\cap C^2(\Omega_T)$ 中的惟一性以及在能量模意义下解对于初值和右端项的连续依赖性。
\end{remark}

\begin{figure}[H]
\centering
\includegraphics[width=0.5\textwidth]{PDE图6.png}
\caption{区域 $\Omega_\tau$ 示意图}
\label{figure:omega_tau}
\end{figure}

\begin{proof}
在波动方程两端同乘 $u_t$，并在区域 $\Omega_\tau$ 上积分，得到
\begin{align*}
\iiint_{\Omega_\tau}u_t[u_{tt}-a^2(u_{xx}+u_{yy})]\,\mathrm{d}x\mathrm{d}y\mathrm{d}t=\iiint_{\Omega_\tau}u_tf\,\mathrm{d}x\mathrm{d}y\mathrm{d}t.
\end{align*}
注意到
\begin{align*}
u_t[u_{tt}-a^2(u_{xx}+u_{yy})]
=\frac{1}{2}[u_t^2+a^2(u_x^2+u_y^2)]_t-a^2[(u_tu_x)_x+(u_tu_y)_y],
\end{align*}
将左端记为 $I$，应用 \hyperref[Basis of Analytics-theorem:Gauss-Green公式(散度定理)]{Gauss-Green公式}得
\begin{align*}
I=\iint_{\partial\Omega_\tau}\left\{\frac{1}{2}[u_t^2+a^2(u_x^2+u_y^2)]\cos(n,t)-a^2[u_tu_x\cos(n,x)+u_tu_y\cos(n,y)]\right\}\mathrm{d}S.
\end{align*}
在柱体 $\Omega\times[0,\tau]$ 的侧面 $\partial\Omega\times[0,\tau]$ 上，$\cos(n,t)=0$，且由边值条件 $u|_{\partial\Omega}=0$ 得 $u_t|_{\partial\Omega}=0$，故侧面上的积分为零。因此
\begin{align*}
I=\frac{1}{2}\iint_\Omega [u_t^2(x,y,\tau)+a^2(u_x^2(x,y,\tau)+u_y^2(x,y,\tau))]\,\mathrm{d}x\mathrm{d}y
-\frac{1}{2}\iint_\Omega[\psi^2+a^2(\varphi_x^2+\varphi_y^2)]\,\mathrm{d}x\mathrm{d}y.
\end{align*}
于是得到
\begin{align}
\frac{1}{2}\iint_\Omega [u_t^2+a^2(u_x^2+u_y^2)]\,\mathrm{d}x\mathrm{d}y-\frac{1}{2}\iint_\Omega[\psi^2+a^2(\varphi_x^2+\varphi_y^2)]\,\mathrm{d}x\mathrm{d}y
=\iiint_{\Omega_\tau}u_tf\,\mathrm{d}x\mathrm{d}y\mathrm{d}t. \label{eq:energy-identity}
\end{align}
当 $f\equiv 0$ 时，上式右端为零，即得等号且 $M=1$。

对于一般情形，由不等式 $2ab\leqslant \varepsilon a^2+\frac{1}{\varepsilon}b^2$（$\varepsilon>0$）得
\begin{align*}
\iiint_{\Omega_\tau}u_tf\,\mathrm{d}x\mathrm{d}y\mathrm{d}t
\leqslant \frac{\varepsilon}{2}\iiint_{\Omega_\tau}u_t^2\,\mathrm{d}x\mathrm{d}y\mathrm{d}t+\frac{1}{2\varepsilon}\iiint_{\Omega_\tau}f^2\,\mathrm{d}x\mathrm{d}y\mathrm{d}t.
\end{align*}
代入 \eqref{eq:energy-identity} 得
\begin{align*}
&\iint_\Omega [u_t^2(x,y,\tau)+a^2(u_x^2(x,y,\tau)+u_y^2(x,y,\tau))]\,\mathrm{d}x\mathrm{d}y \\
&\leqslant \iint_\Omega[\psi^2+a^2(\varphi_x^2+\varphi_y^2)]\,\mathrm{d}x\mathrm{d}y
+\varepsilon\iiint_{\Omega_\tau}u_t^2\,\mathrm{d}x\mathrm{d}y\mathrm{d}t+\frac{1}{\varepsilon}\iiint_{\Omega_\tau}f^2\,\mathrm{d}x\mathrm{d}y\mathrm{d}t.
\end{align*}
记
\begin{align*}
G(\tau)=\iiint_{\Omega_\tau}[u_t^2+a^2(u_x^2+u_y^2)]\,\mathrm{d}x\mathrm{d}y\mathrm{d}t,
\end{align*}
则
\begin{align*}
\frac{\mathrm{d}G(\tau)}{\mathrm{d}\tau}=\iint_\Omega [u_t^2(x,y,\tau)+a^2(u_x^2(x,y,\tau)+u_y^2(x,y,\tau))]\,\mathrm{d}x\mathrm{d}y.
\end{align*}
于是得到微分不等式
\begin{align*}
\frac{\mathrm{d}G(\tau)}{\mathrm{d}\tau}\leqslant \varepsilon G(\tau)+F_\varepsilon(\tau),
\end{align*}
其中
\begin{align*}
F_\varepsilon(\tau)=\iint_\Omega[\psi^2+a^2(\varphi_x^2+\varphi_y^2)]\,\mathrm{d}x\mathrm{d}y+\frac{1}{\varepsilon}\iiint_{\Omega_\tau}f^2\,\mathrm{d}x\mathrm{d}y\mathrm{d}t.
\end{align*}
$F_\varepsilon(\tau)$ 是 $\tau$ 的单调递增函数。由 Gronwall 不等式（引理 \eqref{eq:gronwall-diff}）得
\begin{align*}
G(\tau)\leqslant \int_0^\tau e^{\varepsilon(\tau-t)}F_\varepsilon(t)\,\mathrm{d}t\leqslant \frac{1}{\varepsilon}(e^{\varepsilon\tau}-1)F_\varepsilon(\tau),
\end{align*}
进而
\begin{align*}
\frac{\mathrm{d}G(\tau)}{\mathrm{d}\tau}\leqslant e^{\varepsilon\tau}F_\varepsilon(\tau)\leqslant e^{\varepsilon T}F_\varepsilon(\tau).
\end{align*}
取 $\varepsilon=\frac{1}{T}$，得
\begin{align*}
\iint_\Omega [u_t^2+a^2(u_x^2+u_y^2)]\,\mathrm{d}x\mathrm{d}y
\leqslant e(1+T)\left\{\iint_\Omega[\psi^2+a^2(\varphi_x^2+\varphi_y^2)]\,\mathrm{d}x\mathrm{d}y+\iiint_{\Omega_\tau}f^2\,\mathrm{d}x\mathrm{d}y\mathrm{d}t\right\}.
\end{align*}
定理获证.
\end{proof}

\subsection{广义解}

\begin{definition}
设 $Q_T=(0,l)\times(0,T)$，$\overline{Q}_T=[0,l]\times[0,T]$. 若函数 $u\in C^2(Q_T)\cap C^1(\overline{Q}_T)$ 满足
\begin{align}
\begin{cases}
u_{tt}-a^2u_{xx}=0, & (x,t)\in Q_T, \\
u(x,0)=\varphi(x), & x\in[0,l], \\
u_t(x,0)=\psi(x), & x\in[0,l], \\
u(0,t)=u(l,t)=0, & t\in[0,T],
\end{cases} \label{eq:classical-def}
\end{align}
则称 $u$ 为混合问题 \eqref{eq:wave-gen-mixed} 在 $Q_T$ 上的\textbf{古典解}.
\end{definition}

\begin{definition}
称
\begin{align}
\mathcal{D}=\{\zeta\in C^2(\overline{Q}_T)\mid \zeta(x,T)=\zeta_t(x,T)=0,\ \zeta(0,t)=\zeta(l,t)=0\} \label{eq:test-space}
\end{align}
为\textbf{试验函数类}，其中的函数称为\textbf{试验函数}.
\end{definition}

\begin{definition}
设 $u\in C(\overline{Q}_T)$. 若对任意 $\zeta\in\mathcal{D}$，有
\begin{align}
\iint_{Q_T}u(\zeta_{tt}-a^2\zeta_{xx})\,\mathrm{d}x\mathrm{d}t
+\int_0^l\varphi(x)\zeta_t(x,0)\,\mathrm{d}x
-\int_0^l\psi(x)\zeta(x,0)\,\mathrm{d}x=0, \label{eq:weak-form-def}
\end{align}
则称 $u$ 为混合问题 \eqref{eq:wave-gen-mixed}
\begin{align}
\begin{cases}
u_{tt}-a^2u_{xx}=0, & (x,t)\in(0,l)\times(0,+\infty), \\
u(x,0)=\varphi(x), & x\in[0,l], \\
u_t(x,0)=\psi(x), & x\in[0,l], \\
u(0,t)=u(l,t)=0, & t\in[0,+\infty).
\end{cases} \label{eq:wave-gen-mixed}
\end{align}
的\textbf{广义解}.
\end{definition}
\begin{remark}
由分部积分可知，古典解必为广义解.
\end{remark}

\begin{theorem}
混合问题 \eqref{eq:wave-gen-mixed} 
的广义解若存在则必唯一.
\end{theorem}
\begin{proof}
设混合问题 \eqref{eq:wave-gen-mixed} 有两个解 $u_1,u_2$，并记 $w=u_1-u_2$。由于 $u_1,u_2$ 满足积分等式 \eqref{eq:weak-form}，则对于任意的 $\zeta\in\mathcal{D}$，函数 $w$ 满足
\begin{align}
\iint_{Q_T}w(\zeta_{tt}-a^2\zeta_{xx})\,\mathrm{d}x\mathrm{d}t=0. \label{eq:w-weak}
\end{align}
设 $g(x,t)\in C_0^\infty(Q_T)$，考虑定解问题
\begin{align}
\begin{cases}
\zeta_{tt}-a^2\zeta_{xx}=g(x,t), & (x,t)\in(0,l)\times(0,T), \\
\zeta(x,T)=0, & x\in[0,l], \\
\zeta_t(x,T)=0, & x\in[0,l], \\
\zeta(0,t)=\zeta(l,t)=0, & t\in[0,T].
\end{cases} \label{eq:dual-problem}
\end{align}
作变量替换 $\tau=T-t$ 和函数变换 $\overline{\zeta}(x,\tau)=\zeta(x,T-\tau)$，则 $\overline{\zeta}$ 在 $Q_T$ 上满足混合问题
\begin{align*}
\begin{cases}
\overline{\zeta}_{\tau\tau}-a^2\overline{\zeta}_{xx}=\overline{g}(x,\tau), & (x,\tau)\in(0,l)\times(0,T), \\
\overline{\zeta}(x,0)=0, & x\in[0,l], \\
\overline{\zeta}_\tau(x,0)=0, & x\in[0,l], \\
\overline{\zeta}(0,\tau)=\overline{\zeta}(l,\tau)=0, & \tau\in[0,T],
\end{cases}
\end{align*}
其中 $\overline{g}(x,\tau)=g(x,T-\tau)$。由于 $\overline{g}$ 是光滑函数且在 $\partial Q_T$ 的一个邻域上为 $0$，从而在角点满足相容性条件，由分离变量法知该混合问题存在解 $\overline{\zeta}\in C^2(\overline{Q}_T)$。显然函数 $\zeta(x,t)=\overline{\zeta}(x,T-t)\in\mathcal{D}$ 是定解问题 \eqref{eq:dual-problem} 的解。于是从 \eqref{eq:w-weak} 我们得到，对于任意 $g\in C_0^\infty(Q_T)$，成立
\begin{align}
\iint_{Q_T}wg\,\mathrm{d}x\mathrm{d}t=0. \label{eq:w-g}
\end{align}
下面证明在 $Q_T$ 上 $w\equiv 0$。用反证法。记 $z=(x,t)$。假设存在 $z_0=(x_0,t_0)\in Q_T$ 使得 $w(z_0)\neq 0$。不妨设 $w(z_0)>0$。由连续性，存在以 $z_0$ 为圆心、$r_0$ 为半径的圆盘 $B(z_0,r_0)\subset Q_T$，使得 $w(z)\geqslant \frac{1}{2}w(z_0)$。取 $g\in C_0^\infty(Q_T)$ 使得 $g\geqslant 0$，在 $B(z_0,\frac{r_0}{2})$ 上 $g\equiv 1$，在 $Q_T\setminus B(z_0,r_0)$ 上 $g\equiv 0$。于是
\begin{align*}
\iint_{Q_T}wg\,\mathrm{d}z
=\int_{B(z_0,r_0)}wg\,\mathrm{d}z
\geqslant \int_{B(z_0,\frac{r_0}{2})}wg\,\mathrm{d}z
\geqslant \frac{\pi r_0^2}{8}w(z_0)>0,
\end{align*}
这与 \eqref{eq:w-g} 矛盾，故 $w\equiv 0$。唯一性获证.
\end{proof}

\begin{theorem}
设 $u$ 是混合问题 \eqref{eq:wave-gen-mixed} 的广义解，则存在常数 $M=M(a,T)$ 使得
\begin{align}
\iint_{Q_T}u^2\,\mathrm{d}x\mathrm{d}t
\leqslant M\left(\int_0^l\varphi^2\,\mathrm{d}x+\int_0^l\psi^2\,\mathrm{d}x\right). \label{eq:weak-energy}
\end{align}
\end{theorem}
\begin{proof}
对于广义解 $u$，我们可以找到一列 $u_n\in C_0^\infty(Q_T)$ $(n=1,2,\cdots)$ 使得
\begin{align}
\iint_{Q_T}|u-u_n|^2\,\mathrm{d}x\mathrm{d}t\leqslant \frac{1}{n}. \label{eq:approx}
\end{align}
固定 $n$，考虑定解问题
\begin{align*}
\begin{cases}
\zeta_{tt}-a^2\zeta_{xx}=u_n, & (x,t)\in(0,l)\times(0,T), \\
\zeta(x,T)=0, & x\in[0,l], \\
\zeta_t(x,T)=0, & x\in[0,l], \\
\zeta(0,t)=\zeta(l,t)=0, & t\in[0,T].
\end{cases}
\end{align*}
由分离变量法知该问题存在解 $\zeta_n\in C^2(\overline{Q}_T)$。于是我们在积分等式 \eqref{eq:weak-form} 取试验函数 $\zeta_n$ 得到
\begin{align}
\iint_{Q_T}u\,u_n\,\mathrm{d}x\mathrm{d}t+\int_0^l\varphi(x)\zeta_{nt}(x,0)\,\mathrm{d}x-\int_0^l\psi(x)\zeta_n(x,0)\,\mathrm{d}x=0. \label{eq:weak-n}
\end{align}
由于 $\zeta_n$ 是古典解，由能量不等式可得
\begin{align*}
\int_0^l|\zeta_{nt}(x,0)|^2\,\mathrm{d}x &\leqslant C_1\iint_{Q_T}u_n^2\,\mathrm{d}x\mathrm{d}t,\\
\int_0^l|\zeta_n(x,0)|^2\,\mathrm{d}x &\leqslant C_2\iint_{Q_T}u_n^2\,\mathrm{d}x\mathrm{d}t.
\end{align*}
于是从 \eqref{eq:weak-n} 和不等式 $2ab\leqslant \varepsilon a^2+\frac{1}{\varepsilon}b^2$ 得
\begin{align*}
\iint_{Q_T}u u_n\,\mathrm{d}x\mathrm{d}t
&\leqslant \left(\int_0^l\varphi^2\,\mathrm{d}x\right)^{1/2}\left(\int_0^l|\zeta_{nt}|^2\,\mathrm{d}x\right)^{1/2}
+\left(\int_0^l\psi^2\,\mathrm{d}x\right)^{1/2}\left(\int_0^l|\zeta_n|^2\,\mathrm{d}x\right)^{1/2} \\
&\leqslant \frac{1}{2}\iint_{Q_T}u_n^2\,\mathrm{d}x\mathrm{d}t+C_3\int_0^l[\varphi^2+\psi^2]\,\mathrm{d}x.
\end{align*}
令 $n\to\infty$，利用 \eqref{eq:approx} 得
\begin{align*}
\iint_{Q_T}u^2\,\mathrm{d}x\mathrm{d}t
\leqslant \frac{1}{2}\iint_{Q_T}u^2\,\mathrm{d}x\mathrm{d}t+C_3\int_0^l[\varphi^2+\psi^2]\,\mathrm{d}x,
\end{align*}
移项即得 \eqref{eq:weak-energy}.
\end{proof}

\begin{corollary}
若 $u_i$（$i=1,2$）分别是对应初值 $(\varphi_i,\psi_i)$ 的广义解，则
\begin{align}
\iint_{Q_T}|u_1-u_2|^2
\leqslant M\left(\int_0^l|\varphi_1-\varphi_2|^2+\int_0^l|\psi_1-\psi_2|^2\right). \label{eq:weak-stab}
\end{align}
即广义解在 $L^2$ 意义下连续依赖于初值.
\end{corollary}
\begin{remark}
这说明广义解在 $L^2$ 模意义下对于初值连续依赖。
\end{remark}

\begin{theorem}
设 $\varphi\in C([0,l])$，$\varphi(0)=\varphi(l)=0$，$\varphi'$ 和 $\psi$ 在 $[0,l]$ 上除去有限个第一类间断点外连续，则混合问题 \eqref{eq:wave-gen-mixed} 存在唯一的广义解，且可表示为级数
\begin{align}
u(x,t)=\sum_{n=1}^{\infty}\left(A_n\cos\frac{na\pi}{l}t+B_n\sin\frac{na\pi}{l}t\right)\sin\frac{n\pi}{l}x, \label{eq:gen-sol-series}
\end{align}
其中
\begin{align}
A_n=\frac{2}{l}\int_0^l\varphi(x)\sin\frac{n\pi}{l}x\,\mathrm{d}x,\qquad
B_n=\frac{2}{na\pi}\int_0^l\psi(x)\sin\frac{n\pi}{l}x\,\mathrm{d}x. \label{eq:gen-coeff}
\end{align}
\end{theorem}
\begin{remark}
定理和推论实际上说明定解问题 \eqref{eq:wave-gen-mixed} 在广义解的框架下仍然是适定的。
\end{remark}
\begin{proof}
只需证明存在性。根据关于 $\varphi,\psi$ 的假设，它们可以按特征函数系 $\{\sin\frac{n\pi}{l}x\}$ 展开为
\begin{align*}
\varphi(x)=\sum_{n=1}^{\infty}\varphi_n\sin\frac{n\pi}{l}x,\qquad
\psi(x)=\sum_{n=1}^{\infty}\psi_n\sin\frac{n\pi}{l}x,
\end{align*}
其中 $\varphi_n,\psi_n$ 是 Fourier 系数。由 $\varphi$ 满足的条件，可得 $\varphi_n=\frac{l}{n\pi}\varphi'_n$，其中 $\varphi'_n$ 是 $\varphi'(x)$ 按 $\{\cos\frac{n\pi}{l}x\}$ 展开的 Fourier 系数。记
\begin{align*}
S_N(x)=\sum_{n=0}^{N}\varphi'_n\cos\frac{n\pi}{l}x,\qquad
T_N(x)=\sum_{n=1}^{N}\psi_n\sin\frac{n\pi}{l}x.
\end{align*}
由特征函数系的完备性，有
\begin{align}
\lim_{N\to\infty}\int_0^l|\varphi'(x)-S_N(x)|^2\,\mathrm{d}x=0, \label{eq:phi-approx}\\
\lim_{N\to\infty}\int_0^l|\psi(x)-T_N(x)|^2\,\mathrm{d}x=0. \label{eq:psi-approx}
\end{align}
由 Bessel 不等式可得
\begin{align}
\sum_{n=1}^{\infty}(\varphi'_n)^2\leqslant \int_0^l|\varphi'(x)|^2\,\mathrm{d}x,\qquad
\sum_{n=1}^{\infty}\psi_n^2\leqslant \int_0^l|\psi(x)|^2\,\mathrm{d}x. \label{eq:bessel}
\end{align}
考虑混合问题
\begin{align*}
\begin{cases}
u_{Ntt}-a^2u_{Nxx}=0, & (x,t)\in Q_T, \\
u_N(x,0)=S_N(x), & x\in[0,l], \\
u_{Nt}(x,0)=T_N(x), & x\in[0,l], \\
u_N(0,t)=u_N(l,t)=0, & t\in[0,T],
\end{cases}
\end{align*}
由古典解的存在性定理，其解可表为
\begin{align*}
u_N(x,t)=\sum_{n=1}^{N}\left(A_n\cos\frac{na\pi}{l}t+B_n\sin\frac{na\pi}{l}t\right)\sin\frac{n\pi}{l}x,
\end{align*}
其中 $A_n=\frac{l}{n\pi}\varphi'_n$，$B_n=\frac{l}{n\pi}\psi_n$。显然 $u_N$ 是级数 \eqref{eq:gen-sol-series} 的前 $N$ 项部分和。在 $u_N$ 的方程两边乘以 $\zeta\in\mathcal{D}$，分部积分得
\begin{align}
\iint_{Q_T}u_N(\zeta_{tt}-a^2\zeta_{xx})\,\mathrm{d}x\mathrm{d}t+\int_0^l S_N(x)\zeta_t(x,0)\,\mathrm{d}x-\int_0^l T_N(x)\zeta(x,0)\,\mathrm{d}x=0. \label{eq:weak-N}
\end{align}
令 $u$ 为级数 \eqref{eq:gen-sol-series} 的和函数。先证 $u_N$ 一致收敛到 $u$。记
\begin{align*}
r_N=u-u_N=\sum_{n=N+1}^{\infty}\left(A_n\cos\frac{na\pi}{l}t+B_n\sin\frac{na\pi}{l}t\right)\sin\frac{n\pi}{l}x.
\end{align*}
则
\begin{align*}
|r_N|&\leqslant \sum_{n=N+1}^{\infty}(|A_n|+|B_n|)
=\sum_{n=N+1}^{\infty}\frac{l}{n\pi}\left(|\varphi'_n|+\frac{1}{a}|\psi_n|\right) \\
&\leqslant \sum_{n=N+1}^{\infty}\left(\frac{l}{n\pi}\right)^2+\frac{1}{2}\sum_{n=N+1}^{\infty}\left(|\varphi'_n|^2+\frac{1}{a^2}|\psi_n|^2\right),
\end{align*}
由 \eqref{eq:bessel} 知当 $N\to\infty$ 时右端趋于 $0$，故 $u_N$ 一致收敛到 $u$。由于 $u_N\in C(\overline{Q}_T)$，故 $u\in C(\overline{Q}_T)$。在 \eqref{eq:weak-N} 中令 $N\to\infty$，利用 \eqref{eq:phi-approx} 和 \eqref{eq:psi-approx} 可得
\begin{align*}
\iint_{Q_T}u(\zeta_{tt}-a^2\zeta_{xx})\,\mathrm{d}x\mathrm{d}t+\int_0^l\varphi(x)\zeta_t(x,0)\,\mathrm{d}x-\int_0^l\psi(x)\zeta(x,0)\,\mathrm{d}x=0,
\end{align*}
即 $u$ 满足广义解定义。存在性得证。唯一性已由定理给出.
\end{proof}




\end{document}